\section{Literature Review}
\label{sec:literature}

\subsection{Credit Scoring: Classical Models, Gradient Boosting, and Deep Learning}
Financial and credit-scoring studies consistently identify borrower credit bureau ratings, loan-to-value (LTV) leverage, vehicle attributes, and lending-channel heterogeneity as core determinants of auto loan default~\cite{agarwal2008}. In traditional retail credit benchmarks, Baesens et al.~\cite{baesens2003} showed that well-regularized Logistic Regression (LR) provides a transparent and competitive baseline relative to single decision trees and early neural networks. Updating these benchmarks across diverse credit portfolios, Lessmann et al.~\cite{lessmann2015} demonstrated that tree ensembles---specifically Random Forests~\cite{breiman2001} and Gradient Boosting---outperform standalone linear and tree classifiers with statistical significance, while emphasizing that incremental AUC gains should always be weighed against interpretability and operational cost.

On structured tabular datasets, Gradient Boosted Decision Tree (GBDT) frameworks lead modern practice, notably \textbf{XGBoost}~\cite{chen2016} (regularized second-order boosting with sparsity-aware branching for missing values and dense zeros) and \textbf{LightGBM}~\cite{ke2017} (fast histogram-based leaf-wise tree growth). Whether deep neural networks surpass GBDTs on consumer credit tables remains widely studied. Both Gunnarsson et al.~\cite{gunnarsson2021} and Grinsztajn et al.~\cite{grinsztajn2022} concluded across large empirical benchmarks that GBDTs match or outperform Multilayer Perceptrons (MLPs) and tabular deep networks at much lower tuning cost, because axis-aligned tree splits naturally model sharp thresholds and resist uninformative features. Notably, a recent paper on the LTFS dataset by Wu~\cite{wu2025} reported an unusually large jump from LR ($\text{AUC}=0.62$) to XGBoost with SMOTE ($\text{AUC}=0.89$), yet relied on random train-test splitting without temporal validation, calibration auditing, or fold-isolated resampling controls.

\subsection{Class Imbalance, Calibration, and Cost-Sensitive Decisions}
While oversampling methods such as SMOTE~\cite{chawla2002} are frequently applied by default to imbalanced data, empirical risk-modeling evidence cautions against uncritical resampling. van den Goorbergh et al.~\cite{goorbergh2022} proved that random undersampling and SMOTE do not improve ROC-AUC separation on moderately imbalanced risk data while severely distorting predicted probabilities (ruining calibration). Furthermore, Saito and Rehmsmeier~\cite{saito2015} demonstrated that Precision--Recall curves (PR-AUC) are more informative than ROC curves alone under class imbalance. Decision theory establishes that when a model predicts well-calibrated risk scores on the natural class prior, asymmetric financial trade-offs should be handled through the Bayes-optimal threshold $c_{\text{FP}}/(c_{\text{FP}}+c_{\text{FN}})$~\cite{elkan2001} or principal-weighted loan cost metrics~\cite{bahnsen2014} rather than synthetic data generation.

\subsection{Explainable AI in Credit Risk}
Because consumer lending requires transparent adverse-action explanations, model interpretability is essential alongside predictive accuracy. Building on the game-theoretic SHAP framework~\cite{lundberg2017}, Lundberg et al.~\cite{lundberg2020} introduced \textbf{TreeSHAP} to compute exact Shapley feature attributions and pairwise interactions for tree ensembles, resolving the inconsistency of traditional split-gain counts while cautioning that strongly correlated predictors share credit across tree splits. Extending SHAP to credit governance, Bussmann et al.~\cite{bussmann2021} showed how Shapley auditing and subgroup analysis uncover distinct borrower risk profiles and highlight fairness risks.

Table~\ref{tab:lit_summary} summarizes the main literature themes and how they directly inform our project pipeline.

\begin{table}[htbp]
  \centering
  \caption{Summary of key literature themes and direct implications for our LTFS pipeline.}
  \label{tab:lit_summary}
  \footnotesize
  \begin{tabularx}{\textwidth}{@{}L{3.4cm} L{2.8cm} Y Y@{}}
    \toprule
    \textbf{Theme} & \textbf{Key References} & \textbf{Established Findings} & \textbf{Design Decision Here} \\
    \midrule
    Auto loans \& tabular model choice & \cite{agarwal2008,baesens2003,lessmann2015,gunnarsson2021,grinsztajn2022} & Credit score, LTV, and dealer channels drive risk; GBDTs outperform LR and tabular MLPs with lower compute. & Benchmark LR, DT, RF, XGBoost, LightGBM, and MLP under clean features and OOT testing. \\[3pt]
    GBDT categoricals \& LTFS gaps & \cite{chen2016,ke2017,wu2025,miccibarreca2001} & Zero-inflated tables suit GBDTs; naive target encoding or random splits cause optimism bias. & Use Out-of-Time validation and strict Out-of-Fold Empirical-Bayes target encoding on high-cardinality IDs. \\[3pt]
    Imbalance, calibration \& cost scoring & \cite{chawla2002,goorbergh2022,saito2015,elkan2001,bahnsen2014} & Resampling skews probability calibration without lifting AUC; PR-AUC and cost-based cutoffs reflect risk better. & Compare fold-isolated SMOTE, RUS, and class weights vs.\ natural prior with Bayesian and principal-weighted thresholds. \\[3pt]
    Credit explainability & \cite{lundberg2017,lundberg2020,bussmann2021} & TreeSHAP gives exact local/global explanations; collinear features share attribution. & Apply exact TreeSHAP with Dev-to-OOT stability checks, error cases, and thin-file audits. \\
    \bottomrule
  \end{tabularx}
\end{table}
